\cotefile{94}
\scanpage{14}\mspage{- II -}
(corps de fonctions algébriques) et un rapport sur la Géométrie algébrique

A ce propos le rédacteur signale que , s'il n'a pas eu lui-même de trop grandes difficultés à donner un traitement presque \textquotedbl{}\unsure{gleichberechtigt}\textquotedbl{} des $p$-adiques et séries formelles , son collègue global a dû séparer nombres et fonctions , et ceci pour des raisons qui semblent \deleted{\unsure{valables}} sérieuses Ceci ne peut qu'apporter de l'eau au moulin de ceux qui veulent traiter les corps de fonctions au moyen de courbes , produits de courbes et jacobiennes . Cependant le rédacteur \unsure{demande} qu'on ne s'engage pas à la légère à renoncer à l'analogie entre fonctions et nombres .

Le rédacteur signale les détails suivants :

a) La prop.3 du \S~1 pourrait être améliorée par prolongement de la valeur absolue d'une algèbre normée réelle à sa complexifiée .

b) Les résultats sur les extensions non ramifiées (\S~2,fin du n\textsuperscript{o}3) méritent des énoncés en forme .

c) Le n\textsuperscript{o} des représentants multiplicatifs (\S~2,n\textsuperscript{o}5) est inutile pour la suite Celui de l'exponentielle \deleted{\illegible{1}} (\S~2,n\textsuperscript{o}6) l'est presque , n'étaient l'utilisation de la finitude de $(K:K^{n})$ dans le th. d'existence du corps de classes local (\S~6,n\textsuperscript{o}7) , et le besoin de la valeur de cet indice qu'a eu le collègue global. Par contre l'exponentielle est inutile pour l'inégalité $(K:N(L))\quad L:K$ ($K$ localement compact , $L$ galoisienne) , qui est fondamentale pour le corps de classes : on peut démontrer celle-ci , soit par une astuce krasnérienne (cf.\S~6 n\textsuperscript{o}1) ; soit comme conséquence de la Führerdiskriminantenformel (\S~6,n\textsuperscript{o}8) ; ce n'est pas nettement plus compliqué que la méthode Chevalley .

d) Le n\textsuperscript{o} du groupe de décomposition est mal fichu (\S~3,n\textsuperscript{o}1) . Il faudrait des énoncés en forme au \S~3,n\textsuperscript{o}3 .

e) Le rédacteur est affolé par les amas de cocycles \deleted{\unsure{sur}} du corps de classes local , et ne comprend rien à tous ces calculs \quad Il a l'impression que la méthode des algèbres simples rend les choses plus claires . Mais il sera ravi si le Haut Commissariat à la Topologie algébrique prend les \deleted{\unsure{\S\S|de}} n\textsuperscript{os} 2,3 et 4 du \S~6 à son compte et y met de l'ordre .

\begin{center}\texttt{-{}-{}-{}-{}-{}-{}-{}-{}-{}-{}-{}-{}-{}-{}-{}-{}-{}-{}-}\end{center}
% archive sheet number: 13
\scanpage{15}
% The first draft, including its heading, is enclosed and cancelled by long strokes.
\deleted{\uline{Cor. 11} Supposons maintenant qu'on ait un syst. projectif \inserted{filtrant strict} $(R_i)_{i\in I}$ de Modules coh. sur $U$, pris affine $U=\operatorname{Spec}A$, et \deleted{supposons} soit $F$ le foncteur sur les $\unsure{\mathcal O_U|\mathscr O_U|A_U}$-Modules}
\deleted{\inserted{$F$ --}}

\uline{Corollaire} 11 \unsure{Soit|Soient} $U,\Omega$ comme dessus, \deleted{$U$ pris affine,} \inserted{Supposons} $S=U=\operatorname{Spec}A$, soit $(P_i)$ \unsure{une|un} \deleted{famille} syst. projectif de $A$-modules, et $F$ le foncteur \unsure{qu'il|qui il} proreprésente. Supposons que l'on ait les deux conditions suivantes

(i) $F(\Omega)$ est un $A$-Module de t.f.

(ii) Pour tt $p\in S$ et tt syst. de paramètres $(f_1,\ldots,f_{\unsure{r|n}})$ de $\deleted{\unsure{\mathcal O|\mathscr O}}A_p$, l'application canonique
\[
 F(\Omega)\otimes A_p/(f_1\unsure{,\ldots,|-}f_r)
 \longrightarrow F(\Omega\otimes A_p/(f_1\unsure{,\ldots,|-}f_r,))
\]
est bijective.

Alors le syst. projectif $(P_i)$ est stationnaire.

De façon précise, (i) \deleted{\unsure{signifie|indique}} \inserted{équivaut à dire} que le système inductif des $\operatorname{Hom}(P_i,\Omega)$ est stationnaire. Soit \deleted{\unsure{$P_i$}} $i_0$ tel que $\operatorname{Hom}(P_{i_0},\unsure{\Omega|L})\overset{\sim}{\longrightarrow}\operatorname{Hom}(P_i,\Omega)$ pour $i\geq i_0$. Alors $P_i\simeq P_{i_0}$ pour $i\geq i_0$.

\uline{Dém} Le diagramme

\diagram{\begin{tikzcd}[ampersand replacement=\&,column sep=small,row sep=large,font=\scriptsize]
\underset{i}{\underset{\longrightarrow}{\lim}}\operatorname{Hom}(P_i,\Omega)\otimes A_p/\underline f A_p \arrow[r,"\sim"] \&
\underset{i}{\underset{\longrightarrow}{\lim}}\operatorname{Hom}(P_i,\Omega_p/\underline f\Omega_p)\\
\operatorname{Hom}(P_{i_0},\Omega)\otimes A_p/\unsure{f|\underline f}A_p \arrow[u,"{\unsure{\wr|\downarrow}}"] \arrow[r] \&
\operatorname{Hom}(P_{i_0},\Omega_p/\underline f\Omega_p) \arrow[u,hook',"\alpha"]
\end{tikzcd}}
% archive sheet number: 14
\scanpage{16}
montre que $\alpha$ est un isom, et comme les morphismes de transition de la deuxième $\underset{\longrightarrow}{\lim}$ sont injectifs, il s'ensuit que
\[
 \operatorname{Hom}(P_{i_0},\Omega/\underline f\unsure{U|\Omega})\longrightarrow
 \operatorname{Hom}(P_i,\Omega_p/f\Omega_p)
\]
est bijectif pour $i\geq i_0$, donc $P_{\unsure{i|j}}\longrightarrow P_{i_0}$ localisé en $p$ est un isom (Cor. 10), donc $P_i\longrightarrow P_{i_0}$ est bijectif.

Mettant ensemble prop. 3, lemme 5, \deleted{\unsure{Remarque|lemme 6}} \inserted{Remarque 5 bis}, et le corollaire 10, on trouve

\uline{Théorème} \uline{12} $A$ anneau noeth. quotient d'un régulier, $F$ un foncteur $\operatorname{Mod} A\longrightarrow(\mathrm{Ens})$. Pour que $F$ soit représentable \inserted{par un module de t.f.}, il f et s. qu'il satisfasse les conditions suivantes :

% The curved arrow carries the first line down to the space below the second.
\tikz[overlay,baseline]{\draw[->] (-.1,.09) to[out=210,in=172] (-.1,-.85);}
(ii) $F$ exact à gauche

(\deleted{\unsure{i}}i) $F$ commute aux $\underset{\longrightarrow}{\lim}$ filtrantes

(iii) a) $\forall\,p\in\operatorname{Spec}A$, et tt $M$ de t.f. sur $\widehat A_p$,
\[
 F(M)\overset{\sim}{\longrightarrow}\underset{\longleftarrow}{\lim}\,F(M/\mathfrak m_p^{n+1}M)
\]

b) $\forall\, f\in A$, et tout module de t.f. sur $\widehat A=\underset{\longleftarrow}{\lim}\,A/f^{\unsure{n+1|n}}A$,
\[
 F(M)\overset{\sim}{\longrightarrow}\underset{\longleftarrow}{\lim}\,F(M/f^{n+1}M)
\]

(iv) $\forall\, M$ de t.f. sur $A$, $F(M)$ est de t.f. sur $A$.

(v) La fonction $p\longrightarrow\operatorname{rang}_{k(p)}F(k(p))$ sur $\operatorname{Spec}A$ est constructible

(vi) Pour tt quotient \inserted{\unsure{non nul|normal}} $B$ de $A$, et tout $B$-Module de t.f. $\Omega$, $\exists\,\deleted{\unsure{g}}f\in B$, $f$ non nilpotent, tel que pour tt $p\in\operatorname{Spec}B_f$, et tte suite \deleted{de} \inserted{\deleted{\unsure{régulière}}} paramètres $f_1,\ldots,f_{\unsure{n|m|r}}$ de $B_p$, l'application canonique
% archive sheet number: 15
\scanpage{17}
\[
 F(\Omega)\otimes B_p/\underline f B_p\longrightarrow F(\Omega\otimes B_p/\underline f B_{\unsure{p|q}})
\]
est bijective.

\uline{Dém}, La nécessité claire, celle de (vi) résultant du \deleted{lemme 6, \unsure{du|et}} corollaire 7, et du fait qu'on peut choisir $f$ tel que $B_f$ soit CM (\unsure{donc|alors} $\underline f=(f_1\unsure{-|\ldots}f_{\unsure{n|r}})$ est nécessairement une suite \deleted{régulière} $B_p$-régulière).

\uline{Suffisance}. On \unsure{note|voit} que les conditions envisagées sont stables par passage à un $F_B$ ($B$ quotient de $A$). Cela nous permet \deleted{de} \deleted{supposer, par réc. noeth., que $F_B$ est \unsure{repr.} par $B$-module de t.f., si $B\ne A$.} \margin{(i)(ii) (iii b)(iv)}grâce à prop. \deleted{\unsure{4}}3, de nous borner à trouver un $f\in A$ non nilpotent, tel que $F_f$ soit repr.

\margin{(i)(ii) (iii a)(iv)(v)}Notons que le lemme 5 \inserted{et Remarque 5 bis} impliquent d'autre part que $F$ est \inserted{\unsure{strt|str.}} proreprésentable par un système projectif \inserted{strict} $(\unsure{P|R}_i)$, les $P_i$ de t.f., on est ramené à trouver $f$, non nilpotent, tel que le syst. $P_{if}$ soit stationnaire.

\margin{(vi)}Or cela est assuré par la \deleted{condition} condition (vi), compte tenu du cor. 11.

\begin{center}\uline{cqfd.}\end{center}

\uline{Remarque} 13 Supposons que pour tt quotient intègre $B$ de $A$, $\exists\,f\in B-0$ tel que $B_f$ soit régulier. Cette condition rend \deleted{évidemment} inutile celle que $A$ soit quotient d'un régulier, et entraîne que \unsure{dans|ds} th. 12, $(\mathrm{vi})\Longrightarrow(\mathrm v)$, \deleted{et de (vi), \unsure{=|-}}
% archive sheet number: 16
\scanpage{18}
\uline{Théorème 14} Soient $f:X\longrightarrow S$ propre de prés. finie, $F,G$ Modules de prés. finie sur $X$, avec $G$ $S$-plat. Alors, \deleted{\unsure{pour un}} il existe un \deleted{\unsure{Mo}}Module $P$ sur $S$, de présentation finie, et un isom de foncteurs en le $\mathcal O_S$-Module quasi-cohérent $M$
% slanted
\margin{\uline{NB} Il y a \unsure{un même|même un|un nouveau|le même} th. de représentabilité en \unsure{suppt} pas $F$ de prés. fin., mais seulement quasi-cohérent ; \unsure{mais alors|alors} $P$ est seulement qu. coh., et on se ramène au cas $F$ de prés. finie par passage à la limite.}
\[
 \underline{\operatorname{Hom}}(P,M)\simeq
 f_*\underline{\operatorname{Hom}}(F,G\otimes_{\mathcal O_S} M)
\]
(et par suite, en prenant les $\Gamma$, l'isom fonctoriel en $M$
\[
 \operatorname{Hom}(P,M)\simeq\operatorname{Hom}(F,G\otimes_{\mathcal O_S}M),
\]
qui montre que $P$ est déterminé à isom unique près comme représentant le foncteur en $M$ du deuxième mb). La formation de $P$ commute au changement de base quelconques $S'\longrightarrow S$.

On voit facilement qu'on peut supposer $S=\operatorname{Spec}A$ affine, et alors c'est une question de représentabilité d'un foncteur $\operatorname{Mod} A)\longrightarrow(\mathrm{Ens})$. On est ramené comme d'habitude au cas d'un $A$ de type fini sur $\mathbb Z$, donc quotient d'un régulier. On peut donc faire appel au th. 12.

Les conditions (i) à (iv) sont standard. \deleted{\unsure{Pour la}} \inserted{En vertu} de \unsure{remarque} 13, il reste à vérifier la condition (vi).

\deleted{Condition (v), \unsure{on|il} \unsure{est ramené à} prouver que si $A$ est intègre, $\exists$ \deleted{\unsure{ouvert}} ouvert non vide $U$ de $S=\operatorname{Spec}A$ tel que sur $U$, la fonction $s\longrightarrow\dim\operatorname{Hom}(F_s,G_s)$ \deleted{\unsure{$g_s$}} soit constante. Or on a $\operatorname{Hom}(F_s,G_s)=\Gamma(X_s,\underline{\operatorname{Hom}}(F_s,G_s))$, d'autre part pour presque tout}
On peut supposer $B=A$.
% archive sheet number: 17
\scanpage{19}
En fait, sauf pour la condition (ii) d'exactitude à gauche, nous pouvons vérifier les conditions du th. 12 \deleted{pour} pour un foncteur de la forme
\[
 M\longrightarrow R^{j}f_*(\underline{\operatorname{Ext}}^{i}(F,G\otimes_{\mathcal O_S}M)\unsure{;|,}
\]
\margin{(sans supposer $G$ plat/$S$ ici)}
\[
 [\text{ou aussi}\quad M\longrightarrow\mathbb R^{j}f_*(\mathbb R\unsure{\operatorname{Hom}|\mathcal Hom|\underline{\operatorname{Hom}}}(F,G\otimes_{\mathcal O_S}M)],
\]
i.e. dans le cas affine, pour des foncteurs
\[
 \begin{aligned}
 M&\longrightarrow H^j(X,\underline{\operatorname{Ext}}^i(F,G\otimes_{\mathcal O_S}M))\\
 {}[\text{ou}\quad M&\longrightarrow \operatorname{Ext}^j(X; F,G\otimes_{\mathcal O_S}M)]
 \end{aligned}
\]
La vérification se fait en deux \unsure{crans|coups}.

\uline{Lemme} 15 Soit $f:X\longrightarrow S=\operatorname{Spec}A$, $f$ de t.f. et $A$ noeth, et $F,G$ \deleted{\unsure{comme des}} cohérents sur $X$, $\Omega$ cohérent sur $S$, $N$ entier Alors il existe un ouvert dense $U$ de $S$ tel que pour $p\in U$ et $f_1,\ldots,f_{\unsure{n|r}}$ \deleted{système} \inserted{suite} $\deleted{\unsure{\mathcal O}}A_p$-régulière, on ait
\[
 \underline{\operatorname{Ext}}^i(F,G\otimes_{\mathcal O_S}\Omega)\otimes A_p/\underline f A_p
 \overset{\sim}{\longrightarrow}\underline{\operatorname{Ext}}^i(F,G\otimes_{\mathcal O_S}\Omega\otimes_{\mathcal O_S}A_p/\underline f A_p)
\]
pour $i\leq N$.

% slanted
\margin{cette \unsure{réduction|rédaction} ne sert à rien en fait\ldots{}..}
\uline{Dém} On peut supposer $X$ affine $=\operatorname{Spec}B$. Alors $F$ a une résolution $\unsure{L_{\bullet}|L_{.}|L_{\cdot}}$ par des modules libres de t.f. sur $B$, et
\[
 \begin{aligned}
 \underline{\operatorname{Ext}}^{\unsure{i|j}}(F,\underbrace{G\otimes_{\deleted{\unsure{\mathcal O}}A}\Omega}_{G'})
 &=\deleted{\unsure{H}}H^i(\underline{\operatorname{Hom}}^\bullet(L,G\otimes_{\mathcal O_S}\Omega)\\
 &= H^i(\underbrace{(\check{L}_{\bullet}\otimes_B G)\otimes_{\mathcal O_S}\Omega}_{K^\bullet})
 \end{aligned}
\]
tandis que
\[
 \underline{\operatorname{Ext}}^i(F,G\otimes_{\mathcal O_S}\Omega\otimes_{\deleted{\unsure{\mathcal O}}A}M\deleted{\unsure{fB_p}})=H^i(K^\bullet\otimes M),
\]
et on veut des conditions pour que
\[
 \underline{\operatorname{Ext}}^i_{\mathcal O_X}(F,G')\otimes_{\deleted{\unsure{A}}\mathcal O_S}M
 \overset{\sim}{\longrightarrow}\underline{\operatorname{Ext}}^i_{\mathcal O_X}(F,G'\otimes_A M)
\]
% archive sheet number: 18
\scanpage{20}
i.e.
\[
 H^i(K^\bullet)\otimes_A M\unsure{\simeq|\overset{\sim}{\longrightarrow}} H^i(K\otimes_A M)\qquad\text{si }i\leq N.
\]

Pour ceci, quitte à tronquer $K^\bullet$ en laissant tomber les $K^j$ pour $j\geq N+2$, on peut supposer $K^\bullet$ borné. Il suffit qu'on ait
\[
 \left\{\begin{aligned}
 \operatorname{Tor}^A_i(K^j,M)&=0 &&\text{pour }i>0,\ \text{tout }j\\
 \operatorname{Tor}^A_i(H^j(K^\bullet),M)&=0 &&\text{pour }i>0,\ \text{tout }j.
 \end{aligned}\right.
\]
Comme les modules $K^j,H^j(K)$ sont de t.f. et en nb fini, on est ramené à \deleted{\unsure{voir que|prouver le}} \inserted{prouver le}

\uline{Lemme} 16 \inserted{Soit $R$ cohérent sur $X$.} $\exists\,U$ dense dans $S=\operatorname{Spec}A$ tel que $p\in U$, $f_1\unsure{--|\ldots} f_{\unsure{n|r}}$ \deleted{suite de \unsure{paramètres}} suite $A_p$-régulière implique
\[
 \operatorname{Tor}^A_i(F,A_p/\underline f A_p)=0\qquad\text{pour }i>0.
\]

Or, pour qu'on ait ceci, il suffit, d'après l'argument \unsure{immédiat} du lemme 6, de savoir que $\operatorname{gr}_J(\unsure{\mathcal O_X|B_X})$ \deleted{est} \inserted{et $\operatorname{gr}_J(A)$ soient} plats sur $A/J$ en les points de $U$ ($J=\text{nilradical de } A$). Or on peut choisir $U$ ainsi, par le \uline{th. de platitude générique}.

Le lemme \deleted{\unsure{14|15}}15 montre que \deleted{quitte à \unsure{en bien rappelant|en bien remplaçant} $\unsure{\Omega|B}$, on peut supposer que l'on} \inserted{l'on est} ramené, pour traiter le cas de $\deleted{\unsure{R^j f}}\unsure{R^j|\mathbb R^j|R^q}f_*(\underline{\operatorname{Ext}}^i(F,G\otimes\Omega))$, \deleted{en changeant} à prouver le
% archive sheet number: 19
\scanpage{21}
\uline{Lemme} 17 Soit $f:X\longrightarrow S$ propre, \inserted{$S$ noeth.,} et soit $\underline E$ coh. sur $X$. Alors il existe un ouvert dense $U$ de $S$ tel que $\forall\,s\in S$, et $\underline f=f_1,\ldots f_{\unsure{n|r}}$ suite $\mathcal O_{S,s}$-régulière, l'hom naturel
\[
 R^jf_*(E)_s\otimes\mathcal O_{S,s}/\underline f\mathcal O_{S,s}
 \longrightarrow R^j f^{(s)}_*(E^{(s)}\otimes\deleted{\unsure{A}}\mathcal O_{S,s}/\underline f\mathcal O_{S,s})
 \tag*{$\alpha$)}
\]
soit un isom pour tt $j$.
\margin{($f^{(s)}:X^{(s)}=X\times_S\operatorname{Spec}\mathcal O_{S,s}\longrightarrow\operatorname{Spec}\mathcal O_{S,s}$, $E^{(s)}$ induit par $E$ sur $X^{(s)}$)}

\uline{Dém} Utilisant le lemme 16, on voit qu'on peut supposer déjà que
\[
 E^{(s)}\otimes\mathcal O_{S,s}/\underline f\mathcal O_{S,s}
 \overset{\sim}{\longleftarrow}E^{(s)}\overset{\mathbb L}{\otimes}\underline{O}_{S,s}/\underline f\mathcal O_{S,s}.
\]
De \unsure{m\textsuperscript{e}|même}, on peut supposer \deleted{que} (faisant $X=S$, \unsure{..} et représentant $Rf_*(E)$ par un complexe borné à composantes cohérentes $K^\bullet$ sur $S$) que
\[
 K^\bullet\otimes \underline{O}_{S,s}/\underline f\mathcal O_{S,s}
 =K^\bullet\overset{\mathbb L}{\otimes}\mathcal O_{S,s}/\underline f\mathcal O_{S,s}
\]
de sorte que l'hom $\alpha$, grâce à l'isom de projection (où $M\unsure{:|=}A_p/\underline f A_p$)
\[
 Rf_*(\underbrace{E\overset{\mathbb L}{\otimes}_{\mathcal O_S}M}_{E\otimes_{\mathcal O_S}M})
 \longrightarrow\underbrace{Rf_*(E)}_{K^\bullet}\overset{\mathbb L}{\otimes}_{\underline{O}_S}M,
\]
s'identifie à
\[
 H^j(K^\bullet)\otimes M\unsure{-|\longrightarrow} H^j(\overbrace{K^\bullet\otimes M}^{\unsure{\prime\prime|=}K^\bullet\overset{\mathbb L}{\otimes}M}).
\]
Si on \unsure{rapetisse|restreint} encore $U$, on peut supposer
\[
 \operatorname{Tor}_i^{\mathcal O_S}(H^j(K^\bullet),M)=0\qquad\text{pour }i>0,\ \text{tout }j,
\]
et on sait qu'alors l'hom précédent est un isom pour tt $j$. \quad cqfd.
% archive sheet number: 20
\scanpage{22}
\uline{Remarque 18} Il y a lieu d'appeler un foncteur \inserted{(additif ?)} satisfaisant les conditions du th. 12, sauf l'exactitude à gauche (ii), un \uline{foncteur constructible}, Il faudrait \deleted{(pour $X$ sur $Y$)} \unsure{définir|dégager|disposer} de \unsure{m\textsuperscript{e}|même} la notion de \unsure{foncteur|foncteurs} \unsure{constructible|constructibles} $\operatorname{Qcoh}(Y)\longrightarrow\operatorname{Qcoh}(X)$. Il y aurait un \unsure{sorite|suite} : composé de foncteurs constructibles est \unsure{itou|item}, foncteurs $\underline{\operatorname{Ext}}^i,\underline{\operatorname{Tor}}_j,R^{\unsure{n|q}} f_*$ sont constructibles \ldots{}...

\begin{center}\uline{Applications}\end{center}

\uline{Théorème} 19 Soit $f:X\longrightarrow S$ morphisme propre, de prés. finie et plat. Alors pour tt $Z$ \uuline{affine} sur $X$, $f_*(Z)=\displaystyle\prod_{X/S} Z/X$ est représentable par un schéma affine sur $S$. Si $Z/X$ est de prés. finie (resp de type fini) alors il en est de \unsure{m\textsuperscript{e}|même} de $\displaystyle\prod_{X/S}Z/X$ sur $S$.

\uline{Dém} On peut supposer $S$ affine, \deleted{soit} \inserted{donc $X$ qu. cpt \unsure{et|+} \inserted{(\unsure{qu|qu.})} séparé} \deleted{\unsure{$Z_0\quad Z$}} Alors $Z=\underset{\longleftarrow}{\lim}\,Z_i$ \inserted{($\underset{\longleftarrow}{\lim}$ filtrante)}, $Z_i$ affine de présentation finie sur $X$, d'où $f_*(Z)=\underset{\longleftarrow}{\lim}\,f_*(Z_i)$, et il suffit de prouver pour le premier point que les $f_*(Z_i)$ sont repr. par schémas affines de prés. finie sur $S$. Pour le deuxième point, il suffira alors de traiter le cas $Z$ de t.f. sur $X$.
% archive sheet number: 21
\scanpage{23}
Alors on sait que l'on peut prendre les morphismes \inserted{de transition} $Z_j\unsure{\mathrel{\text{---}}|\longrightarrow} Z_i$ des \unsure{imm|imm.} fermées, et il suffit de prouver qu'alors les $f_*(Z_j)\mathrel{\text{---}} f_*(Z_i)$ sont aussi des imm. fermées, car à la limite on aura que $f_*(Z)\hookrightarrow f_*(Z_i)$ par imm. fermée, donc $f_*(Z)$ sera de t.f. sur $S$.

On est donc ramené à prouver en \unsure{m\textsuperscript{e}|même} temps le

\uline{Corollaire 20}. Si $Z\longrightarrow Z'$ est une imm. fermée de $X$-schémas affines, alors $f_*(Z)\longrightarrow f_*(Z')$ est une imm. fermée.

Pour prouver ceci, on est \unsure{également|aisément}, par dévissage standard, ramené au cas où $Z,Z'$ de présentation finie.

Ayant des situations de présentation finie, on est ramené au cas $S$ de type fini sur $\operatorname{Spec}\mathbb Z$, donc noeth.

\uline{1\textsuperscript{er} Cas} $Z$ est un fibré vectoriel $V(F)$, avec $F$ cohérent sur $X$. Alors le th. 12 nous fournit un $P$ cohérent sur $S$, tel que
\[
 \underline{\operatorname{Hom}}(F,f^*(M))\simeq\underline{\operatorname{Hom}}(P,M),
\]
et on voit tout de suite que $f_*(V(F))\simeq V(P)$.

\uline{2\textsuperscript{\unsure{ième|ème}} Cas} Cas Général. On sait qu'on a
% archive sheet number: 22
\scanpage{24}
une immersion fermée
\[
 Z\hookrightarrow Z'=V(\underline F),
\]
et il suffit de prouver que
\[
 f_*(Z)\longrightarrow f_*(Z')
\]
est représentable par immersions fermées. (En \unsure{m\textsuperscript{e}|même} temps, on aura prouvé le cor. 20 !). On est ramené à la situation où on a

\diagram{\begin{tikzcd}[ampersand replacement=\&,column sep=tiny,row sep=small]
 Z \arrow[dr] \& \subset \& Z' \arrow[dl]\\
 \& X \arrow[ur,bend right=38] \arrow[d,dash] \&\\
 \& S \&
\end{tikzcd}}

une section $s'$ de $Z'$ sur $X$, et on veut représenter le ss foncteur \inserted{de $S$} en $T/S$ qui \deleted{\unsure{rend cette}} ``\unsure{fait|font} passer $s'$ par $Z$''). Soit $s'^{-1}(Z)=Y\subset X$, sous-schéma fermé de $X$. \deleted{On} Le foncteur considéré n'est autre que $f_*(Y)=\displaystyle\prod_{X/S}Y/X$. On sait qu'il est représentable par un sous-schéma fermé de $S$ par l'exposé Murre, th. 3.

\begin{center}cqfd.\end{center}

\uline{Remarque} 21 Soit $f:X\longrightarrow S$ propre, plat, de prés. finie. Pour tout Module qu. coh. $F$ sur $X$, on peut trouver un Module qu. coh. $P$ sur $S$ tel que $f_*(\unsure{\mathbb V|V}(F))\simeq \mathbb V(P)$ : on procède par passage à la limite à partir des $F$ de prés. finie.
% archive sheet number: 23
